  \subsection{Proof of Lemma~\ref{lem:lvar_unique}}
  \label{apdx:lvar_unique}
  \begin{lemma}[Uniqueness of the Labeled Variables]
    Given any program $c \in \cdom$ with its label variable set $\lvar(c)$,
    every two labeled variables in this set are different.
    \[
      \forall c \in \cdom, x^i, y^j \in \lvar(c) \sthat i \neq j.
      \]
  \end{lemma}
  \begin{proof}
    This is proved directly by the consistency property of the command label.
  \end{proof}
  
  \subsection{Proof of Lemma~\ref{lem:tracenondec}}
  \label{apdx:tracenondec}
  \begin{lemma}[Trace Non-Decreasing]
    For any program $c \in \cdom$ and initial trace $\trace_0 \in \ftdom_0(c)$,
    if there exists $\trace \in \tdom$ and $c' \in \cdom $ such that $\config{c, \trace_0} \rightarrow \config{c', \trace} $
    then there exists a trace $\trace' \in \tdom$ such that $\trace_0 \tracecat \trace' = \trace$ formally as follows.
    %
    \[
      \begin{array}{l}
      \forall \trace_0 \in \ftdom_0(c), \trace \in \tdom, c, c' \in \cdom \sthat
      % \Big( 
        \config{c, \trace_0} \rightarrow \config{c', \trace} 
      % \lor  \config{c, \trace_0} \uparrow^{\infty} \trace \Big)
      \\ \quad
      \implies \exists \trace' \in \tdom \sthat \trace_0 \tracecat \trace' = \trace 
      \end{array}
    \]
    \end{lemma}
      \begin{proof}
        Taking arbitrary program $c$, initial trace $\trace_0 \in  \ftdom_0(c)$, let
        $\trace \in \ftdom$ be the execution trace and $c'$ be a program satisfying
        \[
          \config{c, \trace_0} \rightarrow \config{c', \trace},
        \]
      then we have the following cases by induction on $c$.
      \caseL{$c = \clabel{\assign{x}{\expr}}^{l}$}
      By the evaluation rule $\rname{assn}$, we have
      $
      {
      \config{\clabel{\assign{x}{\aexpr}}^{l},  \trace_0 } 
      \xrightarrow{} 
      \config{\eskip, \trace :: ({x}, l, v, \bullet)}
      }$, for some $v \in \mathbb{N}$.
      \\
      Picking $\trace = \trace ::({x}, l, v, \bullet)$ and $\trace' =  [({x}, l, v, \bullet) ]$,
      it is obvious that $\trace_0 \tracecat \trace' = \trace$.
      \\
      This case is proved.
      \caseL{$c = \clabel{\assign{x}{\query(\qexpr)}}^{l}$}
      By the evaluation rule $\rname{query}$
      \[
          \inferrule
          {
            \config{\trace, \qexpr} \qarrow \qval
           \and 
          \query(\qval) = v
          \and 
          \event = ({x}, l, v, \qval)
          }
          {
          \config{{\clabel{\assign{x}{\query(\qexpr)}}^l, \trace}}
          \xrightarrow{} 
          \config{{\clabel{\eskip}^l,  \trace \traceadd \event} }
          },
        \]
      let $v \in \mathbb{N}$ be the result of query request and $\qval$ be the query value.
      Then we have
      \[
      {
        \config{{\clabel{\assign{x}{\query(\qexpr)}}^l, \trace_0}}
        \xrightarrow{} 
        \config{{\clabel{\eskip}^l,  \trace_0 \traceadd [({x}, l, v, \qval)]} }.
      }\]
      Let $\trace = \trace_0 \traceadd [({x}, l, v, \qval)]$ and $\trace' =  [({x}, l, v, \qval) ]$,
      then we have
        \[
        \trace_0 \tracecat \trace' = \trace,
        \]
      and this case is proved.
  %     \caseL{$\ewhile \clabel{\bexpr}^{l_w} \edo (c_w)$}
  %     By the first rule applied to $\ewhile \clabel{\bexpr}^{l_w} \edo (c_w)$, there are two subcases:
  %     \subcaseL{$\rname{while-t}$}
  %     If the first rule applied to is $\rname{while-t}$, we have
  %     \\
  %     $\config{{\ewhile \clabel{\bexpr}^{l_w} \edo (c_w), \trace_0}}
  %       \xrightarrow{} 
  %       \config{{
  %       c_w; \ewhile \clabel{\bexpr}^{l_w} \edo (c_w),
  %       \trace :: (b, l_w, \etrue, \bullet)}}~ (1)
  %     $.
  %     \\
  %     Let $\trace_w' \in \ftdom$ be the trace satisfying following execution:
  %     \\
  %     $
  %     \config{{
  %     c_w,
  %     \trace :: (b, l_w, \etrue, \bullet)}}
  %     \xrightarrow{*} 
  %     \config{{
  %     \eskip, \trace_w'}}
  %   $
  %   \\
  %   By induction hypothesis on sub program $c_w$ with starting trace 
  %   $\trace :: (b, l_w, \etrue, \bullet)$ and ending trace $\trace_w'$, 
  %   we know there exist
  %   $\trace_w \in \ftdom$ such that $\trace_w' = \trace :: (b, l_w, \etrue, \bullet) \tracecat \trace_w$.
  %   \\
  %   Then we have the following execution continued from $(1)$:
  %   \\
  %   $
  %   \config{{\ewhile \clabel{\bexpr}^{l_w} \edo (c_w), \trace}}
  %       \xrightarrow{} 
  %       \config{{
  %       c_w; \ewhile \clabel{\bexpr}^{l_w} \edo (c_w),
  %       \trace :: (b, l_w, \etrue, \bullet)}}
  %       \xrightarrow{*} 
  %       \config{\ewhile \clabel{\bexpr}^{l_w} \edo (c_w), \trace :: (b, l_w, \etrue, \bullet) \tracecat \trace_w}
  %       ~(2)
  %   $
  %   By repeating the execution (1) and (2) until the program is evaluated into $\eskip$,
  %   with trace $\trace_w^{i'} $ for $i = 1, \cdots, n n \geq 1$ in each iteration, we know 
  %   in the $i-th$ iteration,
  %    there exists  $\trace_w^i \in \ftdom$ such that  
  %   $\trace_w^{i'} = \trace_w^{(i-1)'} :: (b, l_w, \etrue, \bullet) ++ \trace_w^{i'}$
  %   \\
  %   Then we have the following execution:
  %   \\
  %   $
  %   \config{{\ewhile \clabel{\bexpr}^{l_w} \edo (c_w), \trace}}
  %       \xrightarrow{} 
  %       \config{{
  %       c_w; \ewhile \clabel{\bexpr}^{l_w} \edo (c_w),
  %       \trace :: (b, l_w, \etrue, \bullet)}}
  %       \xrightarrow{*} 
  %       \config{\ewhile \clabel{\bexpr}^{l_w} \edo (c_w), \trace_w^{n'}}
  %       \xrightarrow{}^\rname{{while-f}}
  %       \config{\eskip, \trace_w^{n'}:: (b, l_w, \efalse, \bullet)}
  %   $ and $\trace_w^{n'} = \trace :: (b, l_w, \etrue, \bullet) \tracecat \trace_w^{1} :: \cdots :: (b, l_w, \etrue, \bullet) \tracecat \trace_w^{n} $.
  %   \\
  %   Picking $\trace' = \trace_w^{n'} :: (b, l_w, \efalse, \bullet)$ and $\trace'' = [(b, l_w, \etrue, \bullet)] \tracecat \trace_w^{1} :: \cdots :: (b, l_w, \etrue, \bullet) \tracecat \trace_w^{n}$,
  %   we have 
  %   $\trace ++ \trace'' = \trace'$.
  %   \\
  %   This case is proved.
  %     \subcaseL{$\textbf{while-f}$}
  %     If the first rule applied to $c$ is $\rname{while-f}$, we have
  %     \\
  %     $
  %     {
  %       \config{{\ewhile \clabel{\bexpr}^{l_w} \edo (c_w), \trace}}
  %       \xrightarrow{}^\rname{while-f}
  %       \config{{
  %       \eskip,
  %       \trace :: (b, l_w, \efalse, \bullet)}}
  %     }$,
  %     In this case, picking $\trace' = \trace ::(b, l_w, \efalse, \bullet)$ and $\trace'' =  [(b, l_w, \efalse, \bullet) ]$,
  %     it is obvious that $\trace \tracecat \trace'' = \trace'$.
  %     \\
  %     This case is proved.
  %     \caseL{$\eif(\clabel{\bexpr}^l, c_t, c_f)$}
  %     This case is proved in the same way as \textbf{case: $c = \ewhile \clabel{\bexpr}^{l} \edo c$}.
  %     \caseL{$c = c_{s1};c_{s2}$}
  %    By the induction hypothesis on $c_{s1}$ and $c_{s2}$ separately,
  %    we have this case proved.
  %   \end{proof}
  %   %
  %   \subsection{Proof of Corollary~\ref{coro:aqintrace}}
  %   \label{apdx:aqintrace}
  %   \begin{corollary}
  %     % \label{coro:aqintrace}
  %     For every event and a trace $\trace \in \ftdom$,
  %     if $\event \in \trace$, 
  %     then there exist another event $\event' \in \eventset$ and traces $\trace_1, \trace_2 \in \ftdom$
  %     such that $\trace_1 \tracecat [\event'] \tracecat \trace_2 = \trace $
  %     with 
  %     $\event$ and $\event'$ equivalent but may differ in their query value.
  %     \[
  %       \forall \event \in \eventset, \trace \in \ftdom \sthat 
  %     \event \in \trace \implies \exists \trace_1, \trace_2 \in \ftdom, 
  %     \event' \in \eventset \sthat (\event \in \trace') \land \trace_1 \tracecat [\event'] \tracecat \trace_2 = \trace  
  %     \]
  %     \end{corollary}
  % \begin{proof}
  %   By unfolding the $\aq \aqin t$, we have the following cases:
  %   %
  %   \caseL{$t = []$} The hypothesis is $\efalse$, this case is proved.
  %   %
  %   \caseL{$t = \aq'::t' \land \aq' \aqeq \aq $}
  %   %
  %   Let $t_1 = []$ and $t_2 = t'$, by unfolding the list concatenation operation, we have:
  %   %
  %   \[
  %       t_1 ++ [\aq'] ++ t_2 = [] ++ [\aq'] ++ t' = t
  %   \]
  %   %
  %   Since $\aq' \aqeq \aq$ by the hypothesis, this case is proved.
  %   %
  %   \caseL{$t = \aq'::t' \land \aq' \aqneq \aq $}
  %   %
  %   By induction hypothesis on $\aq \aqin t'$, we know:
  %   %
  %   \[
  %       \exists t_1', t_2', \aq''. ~ s.t., ~ (\aq \aqeq \aq'') \land t_1' ++ [\aq''] ++ t_2' = t'	
  %   \]
  %   %
  %   Let $t_1 = \aq'::t_1'$ and $t_2 = t_2'$, by unfolding the list concatenation operation, we have:
  %   %
  %   \[
  %       t_1 ++ [\aq''] ++ t_2 = (\aq':: t_1') ++ [\aq''] ++ t_2' = \aq'::t' = t
  %   \]
  %   %
  %   Since $\aq'' \aqeq \aq$ by the hypothesis, this case is proved.    %
  %   
  \caseL{$c = c_{1};c_{2}$}
  By the first rule applied to $c$ in operational semantics, there are two cases as follows.
  \subcaseL{$\rname{seq1}$}
  If the first rule applied to $c$ is $\rname{seq1}$, we have
  \[
    \inferrule
    {
    \config{{c_1, \trace_0}}
    \xrightarrow{}
    \config{{c_1',  \trace_1}}
    }
    {
    \config{{c_1; c_2, \trace_0}} 
    \xrightarrow{} 
    \config{{c_1'; c_2, \trace_1}}
    }
  \]
  By induction hypothesis on $c_1$, we know 
  \[
    \exists \trace_1' \in \tdom \sthat \trace_0 \tracecat \trace_1' = \trace_1
  \]
  Let $\trace = \trace_1$ and $\trace' =  \trace_1'$,
  then we have $\trace_0 \tracecat \trace' = \trace_0 \tracecat \trace_1' = \trace_1 = \trace$ and this case is proved.
  \subcaseL{\rname{seq2}}
  If the first rule applied to $c$ is $\rname{seq2}$, we have
  \[
    \inferrule
    {
    \config{{c_2, \trace_0}}
    \xrightarrow{}
    \config{{c_2',  \trace_2}}
    }
    {
    \config{{\clabel{\eskip}^l; c_2, \trace_0}} 
    \xrightarrow{} 
    \config{{c_2, \trace_2}}
    }
  \]
  By induction hypothesis on $c_2$, we know 
  \[
    \exists \trace_2' \in \tdom \sthat \trace_0 \tracecat \trace_2' = \trace_2
  \]
  Let $\trace = \trace_2$ and $\trace' =  \trace_2'$,
  then we have $\trace_0 \tracecat \trace' = \trace$ and this case is proved.
  \caseL{$\ewhile \clabel{b}^{l_w} \edo (c_w)$}
  By the first rule applied to $c$ in operational semantics, there are two cases as follows.
  \subcaseL{$\rname{while-t}$}
  If the first rule applied to $c$ is \rname{while-t}, we have
  \[
    \config{{\ewhile \clabel{b}^{l_w} \edo (c_w), \trace_0}}
    \xrightarrow{} 
    \config{{
    c_w; \ewhile \clabel{b}^{l_w} \edo (c_w),
    \trace_0 :: (b, l_w, \etrue, \bullet)}}.
  \]
  Let $\trace = \trace_0 ::(b, l_w, \etrue)$ and $\trace' =  [(b, l_w, \etrue)]$,
  then we have $\trace_0 \tracecat \trace' = \trace$ and this case is proved.
  \subcaseL{$\rname{while-f}$}
  If the first rule applied to $c$ is $\rname{while-f}$, we have
  \[
    \config{{\ewhile \clabel{b}^{l_w} \edo (c_w), \trace_0}}
    \xrightarrow{}^\rname{while-f}
    \config{{
    \eskip,
    \trace_0 :: (b, l_w, \efalse, \bullet)}},
  \]
  Let $\trace = \trace_0 ::(b, l_w, \efalse)$ and $\trace' =  [(b, l_w, \efalse)]$,
  then we have $\trace_0 \tracecat \trace' = \trace$ and this case is proved.
  \caseL{$\eif(\clabel{b}^l, c_t, c_f)$}
  By the first rule applied to $\eif(\clabel{b}^l, c_t, c_f)$ in operational semantics,
  we have two cases \rname{if-t} and \rname{{if-f}}.
  Then both of the two cases will append the evaluation result of the guard $\clabel{b}^l$, $(b, l_w, \etrue)$ and $(b, l_w, \efalse)$ respectively.
  \\
  Formally, the two cases will be proved in the same way as \textbf{case: $c = \ewhile \clabel{b}^{l} \edo (c_w)$}.
\end{proof}